\section{Methods and Benchmark Design}
\label{sec:methods}

\subsection{Experimental design and frozen model identity}

% Question: What is the experiment, and what exactly is the method under
% evaluation?
% Answer: A controlled recovery experiment on one method-neutral injected
% detector cube, evaluating NOVA's frozen science release, NOVA-S Science R1,
% against extraction-based analyses of the same cube.
% Evidence: The NOVA-S handoff README and science contract.

The experiment is a controlled recovery test. A detector-level injector
builds a synthetic SOSS time series around a known transmission spectrum,
and NOVA and one or more extraction-based analyses receive the same detector
cube, time grid, uncertainty array, data-quality information, and declared
calibration products. Each method fixes its spectral prediction before the
injected spectrum is revealed, so model development is separated from
absolute evaluation.

NOVA is the general framework introduced in
Section~\ref{sec:structural-limitations-inference}. The implementation
evaluated in this assessment is its frozen science release, NOVA-S
Science R1, called NOVA-S below. The release is immutable and authenticated
by content hashes; later response experiments and injector variants are
candidates or diagnostics, not NOVA-S. One rule governed all development: no
calibration, mask, response, threshold, starting point, or convergence
decision may use the injected truth or any recovered-depth residual.
Everything is built from the observed detector data, instrument reference
files, and out-of-transit (OOT) information, and the injected spectrum is
read only after a prediction has been closed. Model classes that failed
these truth-blind gates are recorded rather than discarded
(Appendix~\ref{app:nova}).

\subsection{The V45 benchmark cube and retained detector data}

% Question: What data does the frozen benchmark contain?
% Answer: 229 Stage-2 rateints integrations whose baseline is the observed
% WASP-17 out-of-transit time series, with a 127-integration synthetic
% event, consumed unchanged by every recovery method.
% Evidence: The exact NOVA-S method note and the injector evolution record.

The frozen benchmark is the V45 cube. Its baseline is the observed
out-of-transit detector time series of the WASP-17 NIRISS/SOSS visit: real
Stage-2 \texttt{rateints} integrations that already contain the target
star, additive background, field sources, and the realised noise. The
injector adds a transit deficit computed from a denoised per-order
expectation of the target source, so only that expectation transits;
background, contaminants, and noise are left untouched. The exact
construction is given in Appendix~\ref{app:benchmarking}. The delivered
cube contains 229 integrations, with the synthetic event occupying
integrations 51--177: 127 event and 102 complement integrations. Stage-2
SCI, ERR, and DQ arrays are consumed directly, and no background
subtraction is applied to the science data before delivery.

The unit of the NOVA-S likelihood is one retained physical detector pixel
in one integration. Valid samples have finite values, positive
uncertainties, and no \texttt{DO\_NOT\_USE} flag, which retains 97{,}557
pixels per integration. The detector support is organised into native
(order, detector-column) groups; six edge groups without reliable mapping
are excluded, leaving 2{,}720. The two orders occupy disjoint support in
this release.

\subsection{The \nova{}-S detector-level forward model}

% Question: What does NOVA-S predict for each detector sample?
% Answer: The sum over orders of a response-propagated product of the
% source-curvature factor, the transit, and native-column continua, plus an
% additive background.
% Evidence: The exact NOVA-S method note, Section 3.

For integration $t$ and physical detector pixel $p$, NOVA-S predicts
\begin{equation}
\begin{split}
 \widehat y_{tp}
 = {}& \sum_o \left[\bm{R}_o\left\{
   g(t)\,\bm{T}_{o,t}(D,\bm{u}_o)\odot
   \bm{\Psi}_o\,\bm{c}_o(t)\right\}\right]_p \\
 & + \left[\bm{H}\bm{\alpha}\right]_{tp} ,
\end{split}
 \label{eq:nova-forward}
\end{equation}
where $o$ indexes the spectral orders and $\odot$ is element-wise
multiplication on the wavelength grid. Reading outward from the star: the
native per-column continua $\bm{c}_o(t)$ set the brightness of each
detector column, $\bm{\Psi}_o$ places those amplitudes on the fine
wavelength grid, $\bm{T}_{o,t}$ dims the stellar light according to the
depth spectrum $D$ and limb darkening $\bm{u}_o$ over each finite exposure,
and $g(t)$ applies a small frozen correction for slow source variation.
The response $\bm{R}_o$ then maps this source model onto detector pixels;
it combines the PASTASOSS wavelength and trace geometry
\citep{BainesEtAl2023Trace,BainesEtAl2023Wavelength}, the WebbKernel
implementation of the ATOCA response mapping
\citep{DarveauBernierEtAl2022,JWSTPipeline2025}, and the empirical q-star
cross-dispersion profile described below. The E13b background
$\bm{H}\bm{\alpha}$ is added last. Both order predictions are summed before
the residual is formed, so the model applies unchanged when order supports
overlap. The background is additive and is never multiplied
by the transit or by $g(t)$, so a change in transit depth cannot be
imitated by dimming the background along with the star.

\subsection{Shared transmission spectrum and transit model}

% Question: How is the transmission spectrum parameterised, and what
% orbital and limb-darkening model generates the transit?
% Answer: 160 shared positive adaptive coefficients split into an achromatic
% depth and a weighted-zero-mean morphology; fixed Keplerian geometry with
% finite-exposure quadrature; eight fitted limb-darkening coordinates around
% the STAGGER theory reference.
% Evidence: The exact NOVA-S method note, Sections 4 and 5.

The transmission spectrum is carried by $K=160$ shared adaptive basis
coefficients,
\begin{equation}
  D_j = D_{\mathrm{global}}\,
  \frac{\exp\!\left[(\bm{H}_v\bm{v})_j\right]}
       {\sum_k w_k \exp\!\left[(\bm{H}_v\bm{v})_k\right]} ,
  \label{eq:depth-basis}
\end{equation}
where $\bm{v}$ is a signed bounded optimiser coordinate, the frozen weights
$w_k$ define the spectral mean, and $\bm{H}_v$ removes the constant gauge.
The achromatic depth $D_{\mathrm{global}}$ and the zero-mean chromatic
morphology are therefore separate quantities, and all depths stay positive.
Orders 1 and 2 evaluate the same physical $D(\lambda)$ on their own
wavelength grids. Two order-discrepancy coordinates exist in the code but
are pinned to zero within $\pm 10^{-12}$ and have no physical freedom. No
atmospheric template, wavelength-local smoother, or truth-informed spectral
prior is used.

The orbital geometry is fixed at the values listed in
Appendix~\ref{app:nova}, with zero eccentricity. Transit light curves come
from the exact Keplerian engine in jaxoplanet \citep{Jaxoplanet2025}, and
each integration is modelled as a finite-exposure average over 21
quadrature points rather than evaluated at a single mid-exposure time.

Quadratic limb darkening \citep{MandelAgol2002} uses the bounded
coordinates of \citet{Kipping2013},
\begin{equation}
 u_1=2\sqrt{q_1}\,q_2, \qquad
 u_2=\sqrt{q_1}\,(1-2q_2).
 \label{eq:kipping-coordinates}
\end{equation}
The wavelength-dependent theory reference is the STAGGER-grid prediction
\citep{MagicEtAl2015}, and the fitted freedom around it is deliberately
small: although six knots per order describe each of $q_1$ and $q_2$, only
a constant offset and a normalised log-wavelength slope are fitted per
variable and order, eight nonlinear coordinates in total. Priors constrain
the deviation from theory and the wavelength curvature, and fine nodes
outside the valid theory support are excluded rather than extrapolated.

\subsection{Response, native continua, uncertainties, and background}

% Question: Which components close the model between transit and detector,
% and how were they calibrated without truth access?
% Answer: An OOT-estimated q-star cross-dispersion response with a
% documented limitation, 5,440 native continuum coefficients, uncertainty
% scales from OOT cross-validation, and the 16-coefficient E13b additive
% background anchored on off-trace pixels.
% Evidence: The exact NOVA-S method note, Sections 6 to 9.

Within each native group, the q-star response describes how the group's
flux is distributed across detector rows. It is a normalised spatial
profile, not an additional spectrum, and it was estimated from the public
V45 cube itself: pixel time series were fitted with a level, a linear
trend, and a depth-free event shape under robust loss, alternating
phase-stratified folds gave independent estimates, and the profile was
admitted only after cross-fold flux and centroid stability gates. The
frozen profile has one documented limitation: its smoother works on an
integer trace-relative coordinate, which produces centroid resets and a
small amplitude mismatch at the reddest columns.
Appendix~\ref{app:nova} gives the construction, the admission gates, and
the defect; its causal assessment belongs to the Results.

Each native group carries a continuum level and a linear time slope,
giving 5{,}440 linear nuisance coefficients that are eliminated exactly at
every nonlinear evaluation. The continuum sets the total brightness of a
group, while q-star sets only its spatial distribution. This split is why
a response-amplitude error relative to the delivered source acts
multiplicatively on the recovered depth.

Detector uncertainties start from the FITS ERR arrays, scaled once per
order, and are then multiplied by a per-detector-row guard derived from
bidirectional OOT cross-validation and floored at one. No wavelength
information, injected truth, or recovered spectrum enters this rule; the
scales and formula are in Appendix~\ref{app:nova}.

The E13b background has eight spatial modes crossed with two temporal
modes, sixteen linear coefficients in all. The spatial family and rank
were chosen by blocked cross-validation on off-trace pixels, and the
second temporal mode was chosen from the leading OOT singular vectors
under a fixed cap on its correlation with the transit reference. The OOT
expectation and normal matrix of these coefficients act as a quadratic
anchor, and background and continua are fitted together in one linear
solve; recovery never subtracts a fixed background. The likelihood carries
no Gaussian-process term for time-correlated residuals: the robust weights
described below handle outliers, and residual temporal correlation is
tracked as a diagnostic.

\subsection{Out-of-transit common source curvature}

% Question: How is slow temporal variation of the source handled?
% Answer: A frozen wavelength-independent curvature factor calibrated on
% held-out OOT folds; the common candidate was selected because the
% per-order candidate failed the requirement that both orders improve.
% Evidence: The exact NOVA-S method note, Section 10.

The factor $g(t)$ corrects for slow variation of the source over the
visit. It was calibrated from the 102 complement integrations, split into
two 51-integration folds: stable trace flux was integrated per order after
subtracting the E13b OOT expectation for this step only, and the ratio of
a weighted quadratic to its affine approximation defined each candidate.
Null, common, and per-order candidates were ranked by held-out
standardised chi-squared. The per-order candidate scored best in aggregate
but failed the requirement that both orders improve individually, so the
common factor was selected. It has unit OOT mean, multiplies only the
stellar and transit bundle, and adds no fitted parameters; its selection
scores and range are in Appendix~\ref{app:nova}.

\subsection{Profiled robust optimisation and convergence}

% Question: How is the model fitted, and what counts as convergence?
% Answer: Exact variable projection of all 5,456 linear coefficients inside
% bounded trust-region least squares, wrapped in exact outer Huber IRLS,
% with preregistered convergence gates, a material-descent audit, and two
% independent deterministic starts.
% Evidence: The exact NOVA-S method note, Section 11, and the frozen code.

Let $\bm{\theta}$ hold the nonlinear spectrum and limb-darkening
parameters and $\bm{\eta}$ the linear continuum and background
coefficients. At fixed $\bm{\theta}$ and fixed robust weights $\bm{W}$,
NOVA-S solves one global linear problem,
\begin{equation}
\begin{split}
 \widehat{\bm{\eta}}(\bm{\theta})=
 \arg\min_{\bm{\eta}}\;
 &\frac{1}{2}\left\|\bm{W}^{1/2}
 [\bm{y}-\bm{A}(\bm{\theta})\bm{\eta}]\right\|_2^2 \\
 &+\frac{1}{2}
 \left\|\bm{L}(\bm{\alpha}-\bm{\alpha}_{\mathrm{off}})\right\|_2^2 ,
\end{split}
 \label{eq:varpro}
\end{equation}
whose regularisation row anchors the background on its OOT expectation
with $\bm{L}^{\mathsf T}\bm{L}$ equal to the OOT normal matrix. Variable
projection \citep{GolubPereyra1973} eliminates all 5{,}440 continuum and
16 background coefficients exactly: the sparse continuum block is
factorised and the background coupling reduces to a 16-dimensional Schur
complement. The profiled Jacobian differentiates through the exact inner
optimum (Appendix~\ref{app:nova}) and is accumulated from streamed
sufficient statistics, so no visit-sized dense Jacobian is ever formed.

The outer nonlinear problem is solved with bounded trust-region reflective
least squares \citep{BranchColemanLi1999} inside exact outer Huber
iteratively reweighted least squares \citep{Huber1964} with threshold
1.345: weights are updated outside the inner solve, so the estimator
remains an exact M-estimator. A fit is accepted only when the robust
weights, the objective, and the shared spectrum have stopped moving, the
scale-aware optimality criteria are met, a deterministic material-descent
audit finds no remaining feasible improvement, the two starts agree, and a
final confirmation fit at the recomputed terminal weights succeeds; the
numeric gates, audit geometry, and cycle caps are in
Appendix~\ref{app:nova}. The two deterministic starts share no state. One begins at the
white-light mean depth recorded in the release configuration, with zero
morphology and theory limb darkening. The other adds a fixed zero-mean
250 ppm morphology perturbation and small fixed limb-darkening offsets. Both starts of the
frozen release converged at robust cycle 14. The accepted solution is
chosen by this convergence policy alone, never by truth error.

\subsection{Benchmark separation, comparison analyses, and fairness status}

% Question: How are injector, recovery, and comparison methods kept
% independent, and what fairness work remains open?
% Answer: Generator, delivery, recovery, and scoring are isolated; one
% opaque cube is consumed unchanged by every pipeline; an authentic Ahsoka
% analysis is the first comparison; the paired group-stage experiment is
% preregistered but not complete.
% Evidence: The injector evolution and fairness record.

Injector and recovery products are separated by construction. Recovery
code reads only the delivered cube and declared auxiliary products; it
cannot see the atmospheric spectrum, generator-only response products, or
scoring code. Hashes identify every delivered array, and the injected
spectrum is loaded only after a prediction and its provenance record are
closed (Appendix~\ref{app:benchmarking}). The injector, in turn, cannot
inspect recovery apertures, q-star products, recovered spectra, or truth
residuals; its source deficit is built independently per order from the
same public instrument references.

The first external comparison is an authentic Ahsoka reduction of the same
delivered cube, following the analysis topology of the published
WASP-17~b spectrum: detector preparation, spectral extraction,
spectroscopic light curves, and channel-level transit inference
\citep{LouieEtAl2025}. Fairness is defined at delivery rather than by
forcing a common internal representation: pixel values, uncertainties,
flags, time grid, and declared wavelength support are identical,
unsupported bins are excluded symmetrically, and each method keeps its own
nuisance model and inference strategy.

One fairness limitation is open. The principal comparison pipelines apply
their most distinctive $1/f$ correction on detector groups before ramp
fitting, and a \texttt{rateints} cube contains no groups. A preregistered
paired experiment therefore injects the same physical deficit at group
stage and at \texttt{rateints} stage and runs NOVA and the comparison
pipelines on both, each with its own truth-blind preprocessing; NOVA is
also run after the same group-level correction. This experiment has not
been completed and is not part of the frozen result. Later benchmark
rounds will run Eureka!, supreme-SPOON/exoTEDRF, transitspectroscopy, and
the official JWST pipeline with ATOCA extraction
\citep{BellEtAl2022,Radica2024,Espinoza2022TransitSpectroscopy,%
JWSTPipeline2025,DarveauBernierEtAl2022,LouieEtAl2025}, each with its
authentic truth-blind settings on the same delivered cube; none of these
pipelines has yet been run with authentic group-level preprocessing on
this benchmark.

\subsection{Reporting support and evaluation metrics}

% Question: On what wavelength support are methods compared, and with which
% metrics?
% Answer: A common R=100 grid restricted prospectively to the direct
% PASTASOSS training domain, scored by absolute RMSE, mean bias, and
% morphology RMSE, with per-recovery diagnostics now and coverage tests in
% the planned ensemble.
% Evidence: The exact NOVA-S method note and appendix draft.

Recovered and injected spectra are compared on a common $R=100$ reporting
grid. One instrument-defined rule fixes the paper-wide support: any bin
whose detector support extends beyond the direct PASTASOSS training domain
is excluded in full, because the wavelength solution beyond the last
training anchor is a polynomial extrapolation. For this visit the rule
keeps 141 bins, the last ending near $2.740\,\mu\mathrm{m}$. The rule is
fixed across methods, changes no fit or prediction, and is applied only
after recovery; its exact detector-coordinate definition is in
Appendix~\ref{app:benchmarking}. Internal regression grids are retained
for diagnostics.

Let $\widehat D_i$ and $D_i^{\mathrm{true}}$ be the recovered and injected
depths on the common grid and $e_i=\widehat D_i-D_i^{\mathrm{true}}$. The
absolute recovery error is
\begin{equation}
 \mathrm{RMSE}_{\mathrm{abs}}=
 \left(\frac{1}{N}\sum_i e_i^2\right)^{1/2},
 \label{eq:absolute-rmse}
\end{equation}
and the morphology error removes the mean residual,
\begin{equation}
 \mathrm{RMSE}_{\mathrm{morph}}=
 \left[\frac{1}{N}\sum_i(e_i-\bar e)^2\right]^{1/2},
 \label{eq:morphology-rmse}
\end{equation}
so that a constant offset is distinguished from an incorrect spectral
shape. Mean and median bias, feature-amplitude error, order-specific
residuals, and subregion morphology are reported alongside. Standardised
residuals and the spectral covariance from the profiled normal system
accompany each individual recovery; empirical coverage of injected depths
requires repeated noise realisations and belongs to the planned ensemble
rather than to the single frozen case. Detector residuals, robust weights,
rank, conditioning, boundary activity, and convergence diagnostics
accompany every score: a low RMSE is not sufficient if a fit fails its
validity gates.
